% Author contribution and AI-use disclosure for module B. Placement: after the last section,
% before the bibliography. On the line paper's final form (paper/ai-statement.tex, the author's
% decision of 2026-09-14, response-plan D5): a short contributor section in the first person --
% roles, verification and review status, the record, responsibility -- with the chronology, the
% figures and their caveats in the repository's process note, here notes/PROCESS-module-b.md
% (exported as notes/PROCESS.md). The hub standard's four rules (writing-style-and-ai-use.md
% section 6) apply: never claim less involvement than there was; numbers from chronicler with
% their date and caveats (in the note, not here); do not repeat the trust boundary, which
% section 1.1 owns; anonymity outranks completeness (not in play).
%
% Written 2026-09-18 from the archivist's report of the same day (archive synced through
% 2026-09-18 20:00 UTC). TO UPDATE BEFORE THE RELEASE: (1) the paragraph on verification and
% review says that no system of another vendor and no human but the author has read the
% article, which is true until the external review rounds have run; rewrite it from their
% record then. (2) The responsibility paragraph does not carry the line paper's clause "and
% have checked every citation": that is a declaration only the author can make, and it is
% left for the author to add if it is true of this article. (3) Re-derive the two figures of
% "The record" when the archive holds the rest of the work; update the note, then this.

\section*{Author contribution and use of AI}

This article was written in close collaboration with generative AI, and the collaboration
went beyond editing assistance at every stage. I write this section in the first person
because it is my own declaration; the article itself is written in the neutral third person.

\paragraph{Roles.} The research programme and the question are mine, as for the line paper
\sscite{fagerstrom2026spatial}. This article is the second of three modules cut from one
working draft, and its mathematics and its formal development were produced in the same
campaign as the line paper's, so that paper's statement and process note cover how they came
about. The decisions of scope, method and policy for this module are mine: the cut between
this module and the next along the dependency graph of the statements; the admission of the
three interface names for the behaviour of a kernel at the origin; that the two constructions
behind the dimension filtration be stated and proved and not only remarked; and the form of
the article, which I set in a review of every section. The article restates what it uses of
the line paper, so that it can be read alone. Its technical sections open with a map. A term
is introduced where it is first used, with its literature. The account of the machine checking
is in \S\ref{sec:machine-checked}, and the body says of a step that it is machine-checked only
inside a few proofs. The numerical run of
\S\ref{sec:numerical-run} was my request.

The prose, the code, the figures and the numerical experiments were written by AI agents
throughout, in dialogue with me; I worked from the conversation and from the built article
rather than from the editor. The Lean~4 development is theirs essentially in full, and none of
the proof terms are mine. So is mathematical content. The admissible cone with its extreme
rays, the bridge from the causal theory and the \Matern{} family as the corner with exact
recursive stages were all named in an agent's reply to my question on the day I asked it,
before any plan existed. The boundedness threshold at the origin and the dimension filtration were proposed by
an agent reading the released line paper in the role of a postdoctoral colleague. The design of
the numerical run is an agent's. During the writing the agents also found and repaired errors
in the statements of record: two while transcribing the statements into the article, and two
in a blind review, a hypothesis that a restating clause had dropped and a uniqueness claim in
the \Matern{} corollary made against too wide a class of realizations. The numerical run then
showed that the text claimed an approximation property for the families with rational stages
which the proposition behind it does not state and which fails. Corrections that make the
article claim less are delegated to the agents by a standing rule of the project, each one
recorded and reported to me, and I keep the veto. The technical decisions of the formalization
were again the agents' recommendations, which I read with their consequences and followed.

\paragraph{Verification and review.} What is machine-checked, and the axioms on which the
verified development rests, are set out in \S\ref{sec:machine-checked}. Review by agents that
had no part in the writing was a working method, as for the line paper: every section was read
blind against the writing standard, one reviewer to a section, and each reviewer restated the
claims from the statements and compared them with the prose, which is how the two errors above
were found. The manuscript was then read by an AI system of another vendor, OpenAI's GPT-6
Astra, in a referee-style review that returned a verdict of major revision and in a
presentation review. Every finding was verified against the text, the sources and the formal development
before it was acted on, and the substantive ones held. The review found what the writing and
the internal reviews had missed: a cited fact on the density at the origin that had been
admitted in a false form, because the record of the source had paraphrased its formula; a
hypothesis missing from the printed statement of the \Matern{} theorem; the realization
corollary stated without its restriction to integer index; a converse the text expected and
that is false; and errors of the numerical example reported under the wrong description. It
also proved a statement the paper had listed as open, which is now clause (4) of
Proposition~\ref{prop:dimension-filtration}; the argument is that system's. A second round
followed the revision. No other vendor's system could be used for it, so it was run as five
referees that were fresh instances of the models that had worked on the article (Claude
Opus~5 and Claude Fable~5.1). Each was given the PDF alone and was forbidden the repository
and the first review, three read the proofs section by section, with the cited theorems read
in the sources, one reproduced the numerical example from its printed description with code
of its own, and one compared the headline claims with the statements and searched the
literature. This round is less independent than the first, since its readers share the blind
spots of the writers, and the false cited fact of the first round was found by the other
vendor's system. It found no false numbered statement. It found an imprecise domain in one
statement, summaries that again claimed more than their statements, a dozen steps that were
true and not on the page, two wrong descriptions in the numerical example and a comparison
made in mixed scales, and prior work that was not cited, the P\'olya frequency functions and
Lindeberg's time-causal limit kernel among it. Several statements were added or widened in
response, all proved in prose. The reviews of both rounds are archived verbatim in the
repository beside the plans that record the disposition of each finding. The statements that
are not machine-checked, which \S\ref{sec:machine-checked} lists and which include the whole
of \S\ref{sec:implementation}, were written by the agents and have been read by me and by
these reviewers, and by no other mathematician. No human reader beyond me has yet reviewed the
article.

\paragraph{The record.} The sessions were archived as the work went on, so the division of
labour can be stated rather than estimated. It is stated, with its dates and its caveats, in
the repository's process note (\texttt{notes/PROCESS.md}), together with the chronology and
the places where the record qualifies a simple account of who did what. The short form: the
work specific to this article took six days on top of the line paper's campaign, and in it
the agents' words outnumber mine by about twenty-four to one. Those figures are lower bounds.
I would not, unaided and as a spare-time activity, have written this article.

\paragraph{Responsibility.} I directed the work throughout and have satisfied myself of the
mathematical content. The cited results were read at their page anchors from held copies by an
agent acting as librarian. Responsibility for this article, including for its errors, is mine
alone; the validity of the machine-checked results depends on the Lean kernel and not on how,
or by whom, the proof terms were produced.
